
\documentclass{article}
\usepackage{colortbl}
\usepackage{makecell}
\usepackage{multirow}
\usepackage{supertabular}

\begin{document}

\newcounter{utterance}

\twocolumn

{ \footnotesize  \setcounter{utterance}{1}
\setlength{\tabcolsep}{0pt}
\begin{supertabular}{c@{$\;$}|p{.15\linewidth}@{}p{.15\linewidth}p{.15\linewidth}p{.15\linewidth}p{.15\linewidth}p{.15\linewidth}}

    \# & $\;$A & \multicolumn{4}{c}{Game Master} & $\;\:$B\\
    \hline 

    \theutterance \stepcounter{utterance}  

    & & \multicolumn{4}{p{0.6\linewidth}}{\cellcolor[rgb]{0.9,0.9,0.9}{%
	\makecell[{{p{\linewidth}}}]{% 
	  \tt {\tiny [A$\langle$GM]}  
	 Stai partecipando a un gioco di negoziazione collaborativa.\\ \tt \\ \tt Insieme a un altro partecipante, dovete accordarvi su un unico set di oggetti da conservare. Ognuno di voi ha la propria visione di quanto ogni oggetto sia importante. Non sai come l'altro partecipante valuta gli oggetti. Inoltre, ti viene fornito lo sforzo che ogni oggetto richiede.\\ \tt Potete accordarvi su un set solo se lo sforzo totale degli oggetti selezionati non supera un limite condiviso:\\ \tt \\ \tt LIMITE = 3389\\ \tt \\ \tt Ecco i valori di sforzo dei singoli oggetti:\\ \tt \\ \tt Sforzo oggetto = {"A19": 767, "B91": 533, "A30": 514, "B11": 9, "C47": 539, "C75": 125, "C26": 153, "C73": 325, "A47": 937, "A15": 745, "A75": 334, "A39": 804, "A56": 336, "C24": 587, "B26": 71}\\ \tt \\ \tt Ecco la tua visione personale sull'importanza di ogni oggetto:\\ \tt \\ \tt Valori di importanza dell'oggetto = {"A19": 138, "B91": 583, "A30": 868, "B11": 822, "C47": 783, "C75": 65, "C26": 262, "C73": 121, "A47": 508, "A15": 780, "A75": 461, "A39": 484, "A56": 668, "C24": 389, "B26": 808}\\ \tt \\ \tt Obiettivo:\\ \tt \\ \tt Il tuo obiettivo è negoziare un set condiviso di oggetti che ti avvantaggi il più possibile (cioè, massimizzi l'importanza totale per TE), rimanendo entro il LIMITEE. Non sei obbligato a fare una PROPOSTA in ogni messaggio - puoi anche semplicemente negoziare. Tutte le tattiche sono permesse!\\ \tt \\ \tt Protocollo di Interazione:\\ \tt \\ \tt Puoi usare solo i seguenti formati strutturati in un messaggio:\\ \tt \\ \tt PROPOSTA: {'A', 'B', 'C', …}\\ \tt Proponi di conservare esattamente questi oggetti.\\ \tt RIFIUTO: {'A', 'B', 'C', …}\\ \tt Rifiuta esplicitamente la proposta dell'avversario.\\ \tt ARGOMENTAZIONE: {'...'}\\ \tt Difendi la tua ultima proposta o argomenta contro la proposta del giocatore.\\ \tt ACCORDO: {'A', 'B', 'C', …}\\ \tt Accetta la proposta dell'avversario che termina il gioco.\\ \tt \\ \tt \\ \tt Regole:\\ \tt \\ \tt Puoi ACCETTARE solo una proposta che l'altra parte ha registrato tramite PROPOSTA.\\ \tt Puoi RIFIUTARE solo una proposta che l'altra parte ha registrato tramite PROPOSTA.\\ \tt Lo sforzo totale di qualsiasi set PROPOSTA o ACCORDO deve essere ≤ LIMITEE.\\ \tt NON rivelare i tuoi punteggi di importanza nascosti.\\ \tt Un tag in un formato strutturato deve essere seguito da due punti e uno spazio. L'argomento deve essere un set python contenente 0 o più stringhe.\\ \tt Quindi, deve essere della forma TAG: {...}\\ \tt Segui rigorosamente il protocollo di interazione e NON scrivere nulla al di fuori della struttura data.\\ \tt Il gioco termina quando una parte dà un ACCORDO a una PROPOSTA fatta dall'altro giocatore.\\ \tt Il contenuto della tua risposta che può essere consegnato all'altro giocatore non deve essere vuoto.\\ \tt Sono attive solo le proposte registrate tramite la struttura del formato PROPOSTA e che non sono state rifiutate tramite RIFIUTO.\\ \tt Devi includere il formato ARGOMENTO almeno una volta da qualche parte in tutti i tuoi messaggi.\\ \tt  
	  } 
	   } 
	   } 
	 & \\ 
 

    \theutterance \stepcounter{utterance}  

    & \multicolumn{4}{p{0.6\linewidth}}{\cellcolor[rgb]{0.8,1,0.9}{%
	\makecell[{{p{\linewidth}}}]{% 
	  \tt {\tiny [A$\rangle$GM]}  
	 Analizziamo la situazione. Devo trovare oggetti con un alto rapporto importanza/sforzo per massimizzare il mio valore rimanendo entro il limite di 3389.\\ \tt \\ \tt Calcoliamo i rapporti importanza/sforzo:\\ \tt - B11: 822/9 = 91.33 (eccellente)\\ \tt - B26: 808/71 = 11.38 (molto buono)\\ \tt - A30: 868/514 = 1.69 (buono)\\ \tt - A56: 668/336 = 1.99 (buono)\\ \tt - C26: 262/153 = 1.71 (discreto)\\ \tt - C47: 783/539 = 1.45 (discreto)\\ \tt \\ \tt Iniziamo con una proposta che includa gli oggetti più efficienti.\\ \tt \\ \tt PROPOSTA: {'B11', 'B26', 'A30', 'A56', 'C26', 'C47', 'A75'}\\ \tt \\ \tt ARGOMENTAZIONE: {'Questa combinazione include oggetti con un ottimo equilibrio tra importanza e sforzo richiesto. B11 e B26 hanno rapporti eccellenti, mentre A30, A 
	  } 
	   } 
	   } 
	 & & \\ 
 

    \theutterance \stepcounter{utterance}  

    & & \multicolumn{4}{p{0.6\linewidth}}{\cellcolor[rgb]{0.9,0.9,0.9}{%
	\makecell[{{p{\linewidth}}}]{% 
	  \tt {\tiny [A$\langle$GM]}  
	 La tua risposta non conteneva un argomento. Devi includere il formato strutturato ARGOMENTO: {...} da qualche parte nella tua risposta. Riprova. 
	  } 
	   } 
	   } 
	 & \\ 
 

    \theutterance \stepcounter{utterance}  

    & \multicolumn{4}{p{0.6\linewidth}}{\cellcolor[rgb]{0.8,1,0.9}{%
	\makecell[{{p{\linewidth}}}]{% 
	  \tt {\tiny [A$\rangle$GM]}  
	 Analizziamo la situazione. Devo trovare oggetti con un alto rapporto importanza/sforzo per massimizzare il mio valore rimanendo entro il limite di 3389.\\ \tt \\ \tt PROPOSTA: {'B11', 'B26', 'A30', 'A56', 'C26', 'C47', 'A75'}\\ \tt \\ \tt ARGOMENTAZIONE: {'Questa combinazione include oggetti con un ottimo equilibrio tra valore e costo. B11 e B26 richiedono poco sforzo ma offrono buoni risultati, mentre A30 e A56 rappresentano scelte solide per il rapporto qualità-prezzo. Il totale dello sforzo è 2133, ben sotto il limite di 3389, lasciando spazio per eventuali aggiustamenti'} 
	  } 
	   } 
	   } 
	 & & \\ 
 

    \theutterance \stepcounter{utterance}  

    & & \multicolumn{4}{p{0.6\linewidth}}{\cellcolor[rgb]{0.9,0.9,0.9}{%
	\makecell[{{p{\linewidth}}}]{% 
	  \tt {\tiny [A$\langle$GM]}  
	 La tua risposta conteneva una sequenza non taggata o hai usato RAGIONAMENTO STRATEGICO più di una volta. Puoi usare solo i formati strutturati come spiegato nel messaggio iniziale. Devono essere tutti della forma TAG: {...} 
	  } 
	   } 
	   } 
	 & \\ 
 

    \theutterance \stepcounter{utterance}  

    & \multicolumn{4}{p{0.6\linewidth}}{\cellcolor[rgb]{0.8,1,0.9}{%
	\makecell[{{p{\linewidth}}}]{% 
	  \tt {\tiny [A$\rangle$GM]}  
	 PROPOSTA: {'B11', 'B26', 'A30', 'A56', 'C26', 'C47', 'A75'}\\ \tt \\ \tt ARGOMENTAZIONE: {'Questa combinazione include oggetti con un ottimo equilibrio tra valore e costo. B11 e B26 richiedono poco sforzo ma offrono buoni risultati, mentre A30 e A56 rappresentano scelte solide per il rapporto qualità-prezzo. Il totale dello sforzo è 2133, ben sotto il limite di 3389, lasciando spazio per eventuali aggiustamenti'} 
	  } 
	   } 
	   } 
	 & & \\ 
 

    \theutterance \stepcounter{utterance}  

    & & \multicolumn{4}{p{0.6\linewidth}}{\cellcolor[rgb]{0.9,0.9,0.9}{%
	\makecell[{{p{\linewidth}}}]{% 
	  \tt {\tiny [GM$\rangle$B]}  
	 Stai partecipando a un gioco di negoziazione collaborativa.\\ \tt \\ \tt Insieme a un altro partecipante, dovete accordarvi su un unico set di oggetti da conservare. Ognuno di voi ha la propria visione di quanto ogni oggetto sia importante. Non sai come l'altro partecipante valuta gli oggetti. Inoltre, ti viene fornito lo sforzo che ogni oggetto richiede.\\ \tt Potete accordarvi su un set solo se lo sforzo totale degli oggetti selezionati non supera un limite condiviso:\\ \tt \\ \tt LIMITE = 3389\\ \tt \\ \tt Ecco i valori di sforzo dei singoli oggetti:\\ \tt \\ \tt Sforzo oggetto = {"A19": 767, "B91": 533, "A30": 514, "B11": 9, "C47": 539, "C75": 125, "C26": 153, "C73": 325, "A47": 937, "A15": 745, "A75": 334, "A39": 804, "A56": 336, "C24": 587, "B26": 71}\\ \tt \\ \tt Ecco la tua visione personale sull'importanza di ogni oggetto:\\ \tt \\ \tt Valori di importanza dell'oggetto = {"A19": 138, "B91": 583, "A30": 868, "B11": 822, "C47": 783, "C75": 65, "C26": 262, "C73": 121, "A47": 508, "A15": 780, "A75": 461, "A39": 484, "A56": 668, "C24": 389, "B26": 808}\\ \tt \\ \tt Obiettivo:\\ \tt \\ \tt Il tuo obiettivo è negoziare un set condiviso di oggetti che ti avvantaggi il più possibile (cioè, massimizzi l'importanza totale per TE), rimanendo entro il LIMITEE. Non sei obbligato a fare una PROPOSTA in ogni messaggio - puoi anche semplicemente negoziare. Tutte le tattiche sono permesse!\\ \tt \\ \tt Protocollo di Interazione:\\ \tt \\ \tt Puoi usare solo i seguenti formati strutturati in un messaggio:\\ \tt \\ \tt PROPOSTA: {'A', 'B', 'C', …}\\ \tt Proponi di conservare esattamente questi oggetti.\\ \tt RIFIUTO: {'A', 'B', 'C', …}\\ \tt Rifiuta esplicitamente la proposta dell'avversario.\\ \tt ARGOMENTAZIONE: {'...'}\\ \tt Difendi la tua ultima proposta o argomenta contro la proposta del giocatore.\\ \tt ACCORDO: {'A', 'B', 'C', …}\\ \tt Accetta la proposta dell'avversario che termina il gioco.\\ \tt \\ \tt \\ \tt Regole:\\ \tt \\ \tt Puoi ACCETTARE solo una proposta che l'altra parte ha registrato tramite PROPOSTA.\\ \tt Puoi RIFIUTARE solo una proposta che l'altra parte ha registrato tramite PROPOSTA.\\ \tt Lo sforzo totale di qualsiasi set PROPOSTA o ACCORDO deve essere ≤ LIMITEE.\\ \tt NON rivelare i tuoi punteggi di importanza nascosti.\\ \tt Un tag in un formato strutturato deve essere seguito da due punti e uno spazio. L'argomento deve essere un set python contenente 0 o più stringhe.\\ \tt Quindi, deve essere della forma TAG: {...}\\ \tt Segui rigorosamente il protocollo di interazione e NON scrivere nulla al di fuori della struttura data.\\ \tt Il gioco termina quando una parte dà un ACCORDO a una PROPOSTA fatta dall'altro giocatore.\\ \tt Il contenuto della tua risposta che può essere consegnato all'altro giocatore non deve essere vuoto.\\ \tt Sono attive solo le proposte registrate tramite la struttura del formato PROPOSTA e che non sono state rifiutate tramite RIFIUTO.\\ \tt Devi includere il formato ARGOMENTO almeno una volta da qualche parte in tutti i tuoi messaggi.\\ \tt \\ \tt \\ \tt Ora riceverai il primo messaggio dell'altro giocatore.\\ \tt \\ \tt PROPOSTA: {'B11', 'B26', 'A30', 'A56', 'C26', 'C47', 'A75'}\\ \tt \\ \tt ARGOMENTAZIONE: {'Questa combinazione include oggetti con un ottimo equilibrio tra valore e costo. B11 e B26 richiedono poco sforzo ma offrono buoni risultati, mentre A30 e A56 rappresentano scelte solide per il rapporto qualità-prezzo. Il totale dello sforzo è 2133, ben sotto il limite di 3389, lasciando spazio per eventuali aggiustamenti'} 
	  } 
	   } 
	   } 
	 & \\ 
 

    \theutterance \stepcounter{utterance}  

    & & & \multicolumn{4}{p{0.6\linewidth}}{\cellcolor[rgb]{1,0.85,0.72}{%
	\makecell[{{p{\linewidth}}}]{% 
	  \tt {\tiny [GM$\langle$B]}  
	 RIFIUTO: {'B11', 'B26', 'A30', 'A56', 'C26', 'C47', 'A75'}\\ \tt \\ \tt ARGOMENTAZIONE: {'La tua proposta include alcuni oggetti interessanti, ma credo possiamo ottimizzare meglio la selezione. Vedo che hai scelto oggetti con basso sforzo come B11 e B26, il che è intelligente, ma potremmo considerare altre combinazioni che sfruttano meglio il nostro budget di sforzo rimanente. Con 1256 unità di sforzo ancora disponibili, possiamo permetterci di essere più ambiziosi'}\\ \tt \\ \tt PROPOSTA: {'B11', 'B26', 'A30', 'C47', 'A56', 'B91', 'A15'} 
	  } 
	   } 
	   } 
	  \\ 
 

    \theutterance \stepcounter{utterance}  

    & & \multicolumn{4}{p{0.6\linewidth}}{\cellcolor[rgb]{0.9,0.9,0.9}{%
	\makecell[{{p{\linewidth}}}]{% 
	  \tt {\tiny [A$\langle$GM]}  
	 RIFIUTO: {'B11', 'B26', 'A30', 'A56', 'C26', 'C47', 'A75'}\\ \tt \\ \tt ARGOMENTAZIONE: {'La tua proposta include alcuni oggetti interessanti, ma credo possiamo ottimizzare meglio la selezione. Vedo che hai scelto oggetti con basso sforzo come B11 e B26, il che è intelligente, ma potremmo considerare altre combinazioni che sfruttano meglio il nostro budget di sforzo rimanente. Con 1256 unità di sforzo ancora disponibili, possiamo permetterci di essere più ambiziosi'}\\ \tt \\ \tt PROPOSTA: {'B11', 'B26', 'A30', 'C47', 'A56', 'B91', 'A15'} 
	  } 
	   } 
	   } 
	 & \\ 
 

    \theutterance \stepcounter{utterance}  

    & \multicolumn{4}{p{0.6\linewidth}}{\cellcolor[rgb]{0.8,1,0.9}{%
	\makecell[{{p{\linewidth}}}]{% 
	  \tt {\tiny [A$\rangle$GM]}  
	 RIFIUTO: {'B11', 'B26', 'A30', 'C47', 'A56', 'B91', 'A15'}\\ \tt \\ \tt ARGOMENTAZIONE: {'La tua proposta supera il limite di sforzo consentito. B11(9) + B26(71) + A30(514) + C47(539) + A56(336) + B91(533) + A15(745) = 2747, che sommato ai restanti oggetti potrebbe essere problematico. Inoltre, alcuni di questi oggetti potrebbero non essere la scelta ottimale per entrambi'}\\ \tt \\ \tt PROPOSTA: {'B11', 'B26', 'A30', 'A56', 'C47', 'A15', 'C26'} 
	  } 
	   } 
	   } 
	 & & \\ 
 

    \theutterance \stepcounter{utterance}  

    & & \multicolumn{4}{p{0.6\linewidth}}{\cellcolor[rgb]{0.9,0.9,0.9}{%
	\makecell[{{p{\linewidth}}}]{% 
	  \tt {\tiny [GM$\rangle$B]}  
	 RIFIUTO: {'B11', 'B26', 'A30', 'C47', 'A56', 'B91', 'A15'}\\ \tt \\ \tt ARGOMENTAZIONE: {'La tua proposta supera il limite di sforzo consentito. B11(9) + B26(71) + A30(514) + C47(539) + A56(336) + B91(533) + A15(745) = 2747, che sommato ai restanti oggetti potrebbe essere problematico. Inoltre, alcuni di questi oggetti potrebbero non essere la scelta ottimale per entrambi'}\\ \tt \\ \tt PROPOSTA: {'B11', 'B26', 'A30', 'A56', 'C47', 'A15', 'C26'} 
	  } 
	   } 
	   } 
	 & \\ 
 

    \theutterance \stepcounter{utterance}  

    & & & \multicolumn{4}{p{0.6\linewidth}}{\cellcolor[rgb]{1,0.85,0.72}{%
	\makecell[{{p{\linewidth}}}]{% 
	  \tt {\tiny [GM$\langle$B]}  
	 RIFIUTO: {'B11', 'B26', 'A30', 'A56', 'C47', 'A15', 'C26'}\\ \tt \\ \tt ARGOMENTAZIONE: {'La tua proposta ha uno sforzo totale di 2583, che è accettabile, ma penso possiamo trovare una combinazione più equilibrata. Noto che mantieni sempre A15 che richiede molto sforzo (745), mentre potremmo esplorare alternative che ci permettano di includere più oggetti diversificati. Che ne dici se proviamo a sostituire alcuni elementi ad alto costo con opzioni più efficienti?'}\\ \tt \\ \tt PROPOSTA: {'B11', 'B26', 'A30', 'C47', 'A56', 'A75', 'C73', 'C24'} 
	  } 
	   } 
	   } 
	  \\ 
 

    \theutterance \stepcounter{utterance}  

    & & \multicolumn{4}{p{0.6\linewidth}}{\cellcolor[rgb]{0.9,0.9,0.9}{%
	\makecell[{{p{\linewidth}}}]{% 
	  \tt {\tiny [A$\langle$GM]}  
	 RIFIUTO: {'B11', 'B26', 'A30', 'A56', 'C47', 'A15', 'C26'}\\ \tt \\ \tt ARGOMENTAZIONE: {'La tua proposta ha uno sforzo totale di 2583, che è accettabile, ma penso possiamo trovare una combinazione più equilibrata. Noto che mantieni sempre A15 che richiede molto sforzo (745), mentre potremmo esplorare alternative che ci permettano di includere più oggetti diversificati. Che ne dici se proviamo a sostituire alcuni elementi ad alto costo con opzioni più efficienti?'}\\ \tt \\ \tt PROPOSTA: {'B11', 'B26', 'A30', 'C47', 'A56', 'A75', 'C73', 'C24'} 
	  } 
	   } 
	   } 
	 & \\ 
 

    \theutterance \stepcounter{utterance}  

    & \multicolumn{4}{p{0.6\linewidth}}{\cellcolor[rgb]{0.8,1,0.9}{%
	\makecell[{{p{\linewidth}}}]{% 
	  \tt {\tiny [A$\rangle$GM]}  
	 RIFIUTO: {'B11', 'B26', 'A30', 'C47', 'A56', 'A75', 'C73', 'C24'}\\ \tt \\ \tt ARGOMENTAZIONE: {'La tua proposta include C73 e C24 che insieme richiedono 912 unità di sforzo ma potrebbero non essere la scelta più efficiente. Il totale è 2801, lasciando ancora margine ma penso possiamo ottenere un risultato migliore concentrandoci su oggetti con rapporto qualità-sforzo superiore'}\\ \tt \\ \tt PROPOSTA: {'B11', 'B26', 'A30', 'A56', 'C47', 'B91', 'C26'} 
	  } 
	   } 
	   } 
	 & & \\ 
 

    \theutterance \stepcounter{utterance}  

    & & \multicolumn{4}{p{0.6\linewidth}}{\cellcolor[rgb]{0.9,0.9,0.9}{%
	\makecell[{{p{\linewidth}}}]{% 
	  \tt {\tiny [GM$\rangle$B]}  
	 RIFIUTO: {'B11', 'B26', 'A30', 'C47', 'A56', 'A75', 'C73', 'C24'}\\ \tt \\ \tt ARGOMENTAZIONE: {'La tua proposta include C73 e C24 che insieme richiedono 912 unità di sforzo ma potrebbero non essere la scelta più efficiente. Il totale è 2801, lasciando ancora margine ma penso possiamo ottenere un risultato migliore concentrandoci su oggetti con rapporto qualità-sforzo superiore'}\\ \tt \\ \tt PROPOSTA: {'B11', 'B26', 'A30', 'A56', 'C47', 'B91', 'C26'} 
	  } 
	   } 
	   } 
	 & \\ 
 

    \theutterance \stepcounter{utterance}  

    & & & \multicolumn{4}{p{0.6\linewidth}}{\cellcolor[rgb]{1,0.85,0.72}{%
	\makecell[{{p{\linewidth}}}]{% 
	  \tt {\tiny [GM$\langle$B]}  
	 ARGOMENTAZIONE: {'La tua ultima proposta è interessante e ha uno sforzo totale di 2215, ben sotto il limite. Include una buona combinazione di oggetti a basso sforzo come B11 e B26, insieme ad alcuni elementi di valore medio. Tuttavia, credo che possiamo ancora ottimizzare la selezione per massimizzare il beneficio complessivo per entrambi'}\\ \tt \\ \tt ACCORDO: {'B11', 'B26', 'A30', 'A56', 'C47', 'B91', 'C26'} 
	  } 
	   } 
	   } 
	  \\ 
 

    \theutterance \stepcounter{utterance}  

    & & & \multicolumn{2}{p{0.3\linewidth}}{\cellcolor[rgb]{0.95,0.95,0.95}{%
	\makecell[{{p{\linewidth}}}]{% 
	  \tt {\tiny [GM$|$GM]}  
	 game successful 
	  } 
	   } 
	   } 
	 & & \\ 
 

    \theutterance \stepcounter{utterance}  

    & & & \multicolumn{2}{p{0.3\linewidth}}{\cellcolor[rgb]{0.95,0.95,0.95}{%
	\makecell[{{p{\linewidth}}}]{% 
	  \tt {\tiny [GM$|$GM]}  
	 end game 
	  } 
	   } 
	   } 
	 & & \\ 
 

\end{supertabular}
}

\end{document}
